%!TEX TS-program = xelatex
%!TEX encoding = UTF-8 Unicode
%%
%% bookspine.tex — 書背（獨立編譯，輸出 bookspine.pdf 後由主文件 \includepdf 匯入）
%%
%% 編譯：xelatex bookspine.tex
%% ════════════════════════════════════════════════════════════════════════
%%  ① 只需修改下方設定區，每個字以「\\ \vfill」或「\\」分隔
%% ════════════════════════════════════════════════════════════════════════

%% -- 系所（左欄）：每字以「\\ \vfill」分隔，依實際字數增減 ----------------
\def\bsDeptChars {%
  航 \\ \vfill
  運 \\ \vfill
  管 \\ \vfill
  理 \\ \vfill
  學 \\ \vfill
  系%
}
\def\bsDeptColH{5.5cm}  %% 欄高：字數 × 約 0.9cm（6字→5.5, 8字→7cm）

%% -- 學校（右欄）：固定 8 字，一般不需更動 ---------------------------------
\def\bsUnivChars {%
  國 \\ \vfill
  立 \\ \vfill
  臺 \\ \vfill
  灣 \\ \vfill
  海 \\ \vfill
  洋 \\ \vfill
  大 \\ \vfill
  學%
}
\def\bsUnivColH{5.5cm}  %% 欄高：8字 × 約 0.9cm

%% -- 其餘各欄：每字以「\\」分隔 --------------------------------------------
\def\bsYear    {2026}                      %% 畢業年份（西元）
\def\bsDegreeV {碩\\士\\學\\位\\論\\文}    %% 學位論文
\def\bsTitleV  {\\數\\理\\統\\計\\數\\理\\統\\計\\數\\理\\統\\計\\數\\理\\統\\計\\數\\理\\統\\計} %% 論文題目（換成實際題目）
\def\bsAuthorV {戴\\忠\\淵}                %% 研究生姓名

%% ════════════════════════════════════════════════════════════════════════
%%  ② 以下排版程式碼無需更動
%% ════════════════════════════════════════════════════════════════════════
\documentclass[a4paper,12pt]{article}

\usepackage{xparse}
\usepackage{fontspec}
\usepackage[slantfont,boldfont]{xeCJK}
\usepackage{graphicx}
\usepackage[a4paper,margin=1cm]{geometry}

%% CJK 字型（fonts/ 目錄優先，否則系統字型）
\IfFileExists{cwTeX Q Kai Medium.ttf}{%
  \setCJKmainfont{cwTeX Q Kai Medium.ttf}[
    Path          = ./,
    Script        = Default,
    BoldFont      = cwTeX Q Kai Medium.ttf,
    ItalicFont    = cwTeX Q Kai Medium.ttf,
    AutoFakeBold  = 2,
    AutoFakeSlant = 0.2
  ]
}{\IfFileExists{fonts/NotoSansTC-Bold.ttf}{%
  \setCJKmainfont{NotoSerifTC-Regular}[
    Path=fonts/, Extension=.ttf,
    UprightFont=*, BoldFont=NotoSerifTC-Bold,
    ItalicFont=*, AutoFakeSlant=0.2
  ]
}{%
  \setCJKmainfont{Noto Serif CJK TC}[
    BoldFont      = Noto Serif CJK TC Bold,
    ItalicFont    = Noto Serif CJK TC,
    AutoFakeSlant = 0.2
  ]
}}

%% \vtext{字\\字\\字} — 垂直堆疊文字
\newcommand{\vtext}[1]{%
  \begin{tabular}[c]{@{}c@{}}#1\end{tabular}%
}

%% ════════════════════════════════════════════════════════════════════════
\begin{document}
\thispagestyle{empty}

\begin{center}
  \setlength{\fboxsep}{0pt}    % 外框貼合書背寬度
  \setlength{\fboxrule}{0.6pt} % 外框粗細

  \framebox{%
    \begin{minipage}[c][270mm]{20mm}
      \centering
      \vspace*{0.8cm}

      %% ── 1. 系所（左）× 學校（右）──────────────────────────────────
      %% minipage [t]=頂對齊  [高度]  [s]=垂直撐滿（vfill 均分字距）
      \begin{tabular}{@{}c@{\hspace{0.2cm}}c@{}}
        \begin{minipage}[t][\bsDeptColH][s]{1.1em}
          \centering\bsDeptChars
        \end{minipage}
        &
        \begin{minipage}[t][\bsUnivColH][s]{1.1em}
          \centering\bsUnivChars
        \end{minipage}
      \end{tabular}

      \vfill

      %% ── 2. 畢業年份（數字橫排，「年」置下）───────────────────────
      \vtext{\bsYear\\年}

      \vfill

      %% ── 3. 學位論文 ──────────────────────────────────────────────
      \vtext{\bsDegreeV}

      \vfill

      %% ── 4. 論文題目 ──────────────────────────────────────────────
      \vtext{\bsTitleV}

      \vfill

      %% ── 5. 研究生姓名 撰 ─────────────────────────────────────────
      \vtext{\bsAuthorV\\撰}

      \vspace*{0.8cm}
    \end{minipage}%
  }
\end{center}

\end{document}
